\PassOptionsToPackage{sols}{handout}
\documentclass[main]{subfiles}
\setcounterrefbeforedot{chapter}{m-ws5-7}
\setcounterrefafterdot{section}{m-ws5-7}
\addtocounter{section}{-1}
\usepackage{solutions}

\begin{document}

\ifSubfilesClassLoaded{%
  \chaptermark{\chapterfive}
}{}

\section{Antiderivatives} \label{ws5-7}

\begin{problemonly}
    \begin{framed}

\textbf{Key Ideas}:

\begin{itemize}[topsep=0pt,partopsep=0pt,parsep=8pt,itemsep=0pt]

\item You should know the difference between definite integrals,
  indefinite integrals, and antiderivatives.

\item You should know indefinite integrals/antiderivatives of the
  following functions: $x^n$ (where $n$ is a constant), $e^x$,
  $b^x$ (where $b$ is a positive constant), $\sin x$, $\cos x$, $\sec^2
  x$, $\sec x \tan x$, $\dfrac{1}{1 + x^2}$, and $\dfrac{1}{\sqrt{1 -
      x^2}}$.

\item You should be able to find antiderivatives of sums, differences, and
  constant multiples of the functions above and be able to use
  these to calculate definite integrals.

\end{itemize}
\end{framed}
\end{problemonly}

\begin{enumerate}
\item 
  \note{This is to reinforce the two-step strategy for such problems:
    first write down an expression involving a definite integral that
    answers the question, then use FTOC 2 to evaluate the definite
    integral.  (Students are often tempted to just write down an
    antiderivative of the rate function and plug in 12.)

    It's also to remind students of FTOC 2 and to motivate the notation
    for indefinite integrals.}

  \begin{problem}[\vfill]
    In January 2010 (which we'll call $t = 0$), the price of a gallon of
    whole milk was \$3.24.  Suppose that, after that, the price changed at
    a rate of $r(t) = 1 + 1.5 \sin t$ cents per month, where $t$ is measured
    in months.  How much did a gallon of whole milk cost 12 months later,
    in January 2011?

    \begin{enumerate}
    \item First, write an expression involving a definite integral that
    answers the question. 
    
    \begin{solution}
        The price 12 months later would be $\boxed{324+
        \displaystyle\int_0^{12} r(t)\,dt}$ cents.
    \end{solution}
    \pspace{\vfill}
    
    \item Use the Fundamental Theorem of Calculus, Part 2 to evaluate
    the integral using an antiderivative.
    
      \begin{solution}
    The cost in cents of a gallon of whole milk at $t = 12$ is 
    \begin{align*}
      324 + \int_0^{12} (1 + 1.5 \sin t)\,dt &= 324 + \left( t - 1.5 \cos t
        \Big|_0^{12} \right) \\
      &= 324 + (12 - 1.5 \cos 12) - (0 - 1.5 \cos 0) \\
      &= 324 + 12 - 1.5 \cos 12 + 1.5 \\
      &= \boxed{337.5 - 1.5 \cos 12} \\
      & \approx 336.23
    \end{align*}
  \end{solution}
    \end{enumerate}
  \end{problem}

  \pspace{\newpage}


% \item 
%   \begin{problem}[\vfill\vfill\newpage]
%     List all of the ``basic'' derivatives that you can. 
%   \end{problem}

%   \begin{solution}
%     You should know the derivatives of power functions $x^n$, exponential
%     functions $b^x$ (including $e^x$), logarithms $\log_b x$ (including
%     $\ln x$), $\sin x$, $\cos x$, $\arcsin x$, $\arccos x$, and $\arctan
%     x$.

%     (You should also know the derivatives of the other trigonometric
%     functions like $\tan x$, but you can figure those out from the
%     derivatives of $\sin x$ and $\cos x$.)
%   \end{solution}

\end{enumerate}
\begin{problemonly}
  \begin{framed}
    We use the notation $\displaystyle \int f(x)\,dx$ to mean ``all
    antiderivatives of $f$'' and call this the \uline{indefinite integral of $f(x)$}.

    While the definite integral $\displaystyle \int_a^b f(x)\,dx$ is a
    number, the indefinite integral $\displaystyle \int f(x)\,dx$ is a
    family of functions.
  \end{framed}
\end{problemonly}

\begin{enumerate}[resume]
\item 
 \textbf{Basic Indefinite Integrals.}
  (\emph{Know these!})

  \begin{itemize}[topsep=0pt]
  \item When $n$ is a constant other than $-1$, $\displaystyle \int x^n\,dx
    = {}$ \solonly{$\displaystyle \bm{\frac{x^{n+1}}{n+1} + C}$}

\probonly{\begin{multicols}{2}}
  \item

    \note{Be sure to talk about what happens for $x < 0$.}

    $\displaystyle \int \frac{1}{x}\,dx = {}$ \solonly{$\displaystyle
      \bm{\ln |x| + C}$}

  \item $\displaystyle \int e^x\,dx = {}$ \solonly{$\displaystyle
      \bm{e^x + C}$} 

    \note{The next one requires trial and error—guess and check and
    guess again. Dignify this as a mathematical problem solving
    strategy.}
\probonly{\end{multicols}}
      
  \item When $b$ is a positive constant, $\displaystyle \int b^x\,dx = {}$
    \solonly{$\displaystyle \bm{\frac{1}{\ln b} b^x + C}$} 

    \probonly{\begin{multicols}{2}}
  \item $\displaystyle \int \sin x\,dx = {}$ \solonly{$\displaystyle
      \bm{- \cos x + C}$} 
  \item $\displaystyle \int \cos x\,dx = {}$ \solonly{$\displaystyle
      \bm{\sin x + C}$} 
  \item $\displaystyle \int \sec^2 x\,dx = {}$ \solonly{$\displaystyle
      \bm{\tan x + C}$} 
  \item $\displaystyle \int \sec x \tan x\,dx = {}$ \solonly{$\displaystyle
      \bm{\sec x + C}$}
  \item $\displaystyle \int \frac{1}{1 + x^2}\,dx = {}$
    \solonly{$\displaystyle \bm{\arctan x + C}$}

  \item

    \note{It's worth talking briefly about why this doesn't violate the
      fact that antiderivatives differ by a constant.}

    $\displaystyle \int \frac{1}{\sqrt{1 - x^2}}\,dx = {}$
    \solonly{$\displaystyle \bm{\arcsin x + C}$ \textbf{or} $\displaystyle
      \bm{-\arccos x + C}$}
      \probonly{\end{multicols}}
  \end{itemize}
\end{enumerate}
\begin{framed}
    Fact: Any 2 antiderivatives of $f$ differ by a constant, but this may not be obvious.
\end{framed}

\begin{enumerate}[resume,topsep=0pt]

\item Match each expression in the first row with the correct
  expression in the second row. \vspace{-8pt}
  
      \null\hfill
      1)\, $\displaystyle \int \frac{1}{1+x^2}\,dx$
      \hfill
      2)\, an antiderivative of $\dfrac{1}{1+x^2}$
      \hfill
      3)\, $\displaystyle \int_0^1 \frac{1}{1+x^2}\,dx$
      \hfill\null
      
        \hrulefill \vspace{-12pt}
        
      \null\hfill 
      a)\, $\arctan x $
      \hfill
      b)\, $\arctan x + C$
      \hfill
      c)\, $1$
      \hfill
      d)\, $\dfrac{\pi}{4}$ 
      \hfill
      e)\, $\dfrac{1}{\tan x} +C$
      \null\hfill

\begin{solution}
\begin{itemize}
    \item 1) matches with b). The indefinite integral 
    $\displaystyle\int f(x)\,dx$ represents
    ``all antiderivatives of $f$'', so we must have $+C$.
    
    \item 2) matches with a). The derivative of $\arctan x$ is 
    $\dfrac{1}{1+x^2}$, so it ONE of its antiderivatives.
    
    \item 3) matches with d). We can use the antiderivative $\arctan x$
    to evaluate the integral,
    \[
    \int_0^1 \frac{1}{1+x^2}\,dx = \arctan(1)-\arctan(0)=\frac{\pi}{4}.
    \]
\end{itemize}
\end{solution}

\pspace{\newpage}


\item  Evaluate the following definite and indefinite integrals.
\begin{framed}
    Useful facts from our previous classes:
    \begin{itemize}[topsep=0pt]
        \item $\displaystyle \int \left[f(x)+g(x)\right]\,dx =
        \int f(x)\,dx + \int g(x)\,dx$
        \item $\displaystyle \int kf(x)\,dx = k\int f(x)\,dx$
    \end{itemize}
\end{framed}

  \begin{enumerate}

  \item

    \begin{problem}[\vfill]
      $\displaystyle \int \left(3e^u -u^2 + 5\right) du$
    \end{problem}

    \begin{solution}
      $\displaystyle \int \left(3e^u - u^2 +5 \right) du = \boxed{3 e^u -
        \frac{u^3}{3} + 5u + C}$.
    \end{solution}

    \item
    \note{A common mistake here is to say $5 \ln |7x| + C$, but they can
      differentiate that to see that they don't get back
      $\frac{5}{7x}$.}

    \begin{problem}[\vfill]
      $\displaystyle \int \frac{5}{7x}\,dx$
    \end{problem}

    \begin{solution}
      First, this is an indefinite integral, so it's really asking for
      all antiderivatives of $\dfrac{5}{7x}$.  Let's rewrite this:
      $\displaystyle \int \frac{5}{7x}\,dx = \int \frac{5}{7} \cdot
      \frac{1}{x}\,dx = \boxed{\frac{5}{7} \ln |x| + C}$.
    \end{solution}

  \item

    \note{This is to emphasize the need for the absolute values in $\ln
      |x|$.  You might get students to visualize this as signed area to
      see whether the answer should be positive or negative.} 

    \begin{problem}[\vfill\newpage]
      $\displaystyle \int_{-e}^{-1/e} \frac{5}{7x}\,dx$
    \end{problem}

    \begin{solution}
      Using our answer to {{{{\#4}}(b)}}, 
      \begin{align*}
        \int_{-e}^{-1/e} \frac{5}{7x}\,dx &= \left. \frac{5}{7} \ln |x|
        \right|_{-e}^{-1/e} \\
        &= \frac{5}{7} \ln \left( \frac{1}{e} \right) - \frac{5}{7} \ln e
        \\ 
        &= \boxed{- \frac{10}{7}}
      \end{align*}
    \end{solution}

  \item
    \begin{problem}[\vfill]
      $\displaystyle \int (e^\pi + ey + \pi^y)\,dy$
    \end{problem}

    \begin{solution}
      $\displaystyle \int (e^\pi + ey + \pi^y)\,dy = \boxed{e^\pi y +
        \frac{e}{2} y^2 + \frac{\pi^y}{\ln \pi} + C}$.  (Keep in mind
      that $e^\pi$ is a constant.)
    \end{solution}

  \item
    \begin{problem}[\stretch{1.5}]
      $\displaystyle \int \left( \frac{e}{x^2} + 2^x \cdot 3^x
      \right)\,dx$
    \end{problem}

    \begin{solution}
      We don't have a rule about integrating products, so we must first
      rewrite $2^x \cdot 3^x$:
      \begin{align*}
        \int \left( \frac{e}{x^2} + 2^x \cdot 3^x \right) dx &= \int
        \left( e x^{-2} + 6^x \right)\,dx \\
        &= - e x^{-1} + \frac{1}{\ln 6} 6^x + C\\
        &= \boxed{- \frac{e}{x} + \frac{1}{\ln 6} 6^x + C}
      \end{align*}
    \end{solution}

  \item

    \note{Some students may want to integrate $\pi + x$ and $\sqrt{x}$
      separately and multiply the results.  Having them differentiate
      their answer is a good way to help them see that this is invalid
      (and that the issue is that they haven't really ``reversed'' the
      Product Rule).}

    \begin{problem}[\stretch{1.5}]
      $\displaystyle \int (\pi+x)\sqrt{x}\, dx$
    \end{problem}

    \begin{solution}
      We don't have a rule about integrating products, so let's first
      rewrite the integrand so that it's no longer a product: 
      \begin{align*}
        \int (\pi + x) \sqrt{x}\,dx &= \int (\pi \sqrt{x} + x
        \sqrt{x})\, dx \\
        &= \int \left( \pi x^{1/2} + x^{3/2} \right) dx \\
        &= \boxed{ \frac{2}{3} \pi x^{3/2} + \frac{2}{5} x^{5/2} + C} 
      \end{align*}
    \end{solution}

  \item

    \note{This is to point out that sometimes, it's still easier to
      interpret as signed area to evaluate integrals.}

    \begin{problem}[\nstretch{1.5}]
      $\displaystyle \int_{-3}^0 \left( \sqrt{9 - x^2} + x^2 \right) dx$
    \end{problem}

    \begin{solution}
      $\displaystyle \int \sqrt{9 - x^2}\,dx$ is not one of the basic
      indefinite integrals that we know from
      {{\#3}}; however, you should recognize
      $\sqrt{9 - x^2}$ as a function whose graph is simple
      geometrically: it's a semicircle.  So, let's split the integral up
      and use two different strategies:
      \[ \int_{-3}^0 \left( \sqrt{9 - x^2} + x \right) dx = \int_{-3}^0
      \sqrt{9 - x^2}\,dx + \int_{-3}^0 x^2 \,dx.\] We can visualize
      $\displaystyle \int_{-3}^0 \left( \sqrt{9 - x^2} + x \right) dx$
      as the area of a quarter circle: 
      \begin{center}
        \begin{tikzpicture}
          \begin{axis}[axis equal=true, xmin=-3.1, xmax=3.1, width=2.5in,
            height=1.5in, xtick={-3, 3}, cycle list=black]
            \addplot+[domain=-3:0, plusarea] {sqrt(9-x^2)} \closedcycle; 
            \addplot+[domain=0:180] ({3*cos(x)}, {3*sin(x)});
          \end{axis}
        \end{tikzpicture}
      \end{center}
      This area is $\dfrac{9 \pi}{4}$.  On the other hand,
      $\displaystyle \int_{-3}^0 x^2 \,dx = \left. \frac{1}{3} x^3
      \right|_{-3}^0 = 9$.  So, $\displaystyle \int_{-3}^0 \left(
        \sqrt{9 - x^2} + x^2 \right) dx = \boxed{\frac{9 \pi}{4} + 9}$.
    \end{solution}

  \end{enumerate}

\item 
  On a camping trip, Ali sets out from his campsite at noon to wander
  along a trail.  We'll describe positions along the trail by how many
  meters north of the campsite they are; for example, position $-10$ is
  $10$ meters south of the campsite.

  Ali's velocity $t$ minutes after noon is $v(t) = 3t^2 - 6t - 9$ meters
  per minute.

  \begin{enumerate}
  \item
    \begin{problem}[\stretch{2}]
      What is Ali's average velocity between noon and 12:04 pm?  Remember
      that average velocity is defined to be $\displaystyle
      \frac{\text{change in position}}{\text{change in time}}$.
    \end{problem}

    \begin{solution}
      The change in Ali's position is
      \begin{align*}
        \int_0^4 v(t)\,dt &=  \int_0^4 (3t^2 - 6t - 9)\, dt \\
        &= t^3 - 3t^2 - 9t \Big|_0^4 \\
        &= -20,
      \end{align*}
      so his average velocity between noon and 12:04 pm is $\dfrac{-20
        \text{ m}}{4 \text{ minutes}} = \boxed{-5 \text{
          meters/minute}}$. 
    \end{solution}

  \item
    \begin{problem}[\stretch{3}]
      What is Ali's average \emph{speed} between noon and 12:04 pm?
      Remember that average speed is defined to be $\displaystyle
      \frac{\text{distance traveled}}{\text{change in time}}$.
    \end{problem}

    \begin{solution}
      Ali's speed function is $|v(t)|$, so the distance he travels between
      noon and 4 pm is $\displaystyle \int_0^4 |v(t)|\,dt$.  To calculate
      this, it might help to first graph $|v(t)|$; $v(t)$ is a quadratic
      that can be factored as $v(t) = 3 (t^2 - 2t - 3) = 3 (t - 3)(t +
      1)$, so the graphs of $v(t)$ and $|v(t)|$ look like this:
      \begin{center}
        \begin{tikzpicture}
          \begin{axis}[xtick={-1, 3}, ytick=\empty, axis on top, tick
            label style={fill=none}]
            \addplot+[domain=-2:5, line width=3pt, skyblue!50!white]
            {3*x^2-6*x-9} node[pos=0.4, below] {$v(t)$}; %
            \addplot+[domain=-2:5, red, thick] { abs(3*x^2-6*x-9) }
            node[pos=0.4, above] {$|v(t)|$};
          \end{axis}
        \end{tikzpicture}
      \end{center}
      In particular, if we focus on the interval $[0, 4]$, we see that
      \[ |v(t)| =
      \begin{cases}
        -v(t) = - 3t^2 + 6t + 9 & \text{ on $[0, 3]$} \\
        v(t) = 3t^2 - 6t - 9 & \text{ on $[3, 4]$}
      \end{cases}\]
      Therefore,
      \begin{align*}
        \int_0^4 |v(t)|\,dt &= \int_0^3 (-3t^2 + 6t + 9)\,dt + \int_3^4
        (3t^2 - 6t - 9)\,dt \\ 
        &= \left( -t^3 + 3t^2 + 9t \Big|_0^3 \right) + \left( t^3 - 3t^2 -
          9t \Big|_3^4 \right) \\
        &= (27 - 0) + \left[ -20 - (-27) \right] \\
        &= 34
      \end{align*}
      Therefore, Ali's average speed is $\dfrac{34 \text{ m}}{4 \text{
          min}} = \boxed{\frac{17}{2} \text{ meters/minute}}$. 
    \end{solution}

  \item

    \begin{problem}[\vfill]
      True or false: The average velocity of a moving object between time
      $a$ and time $b$ is simply the average of its velocities at times
      $a$ and $b$. 
    \end{problem}

    \begin{solution}
      False, and this problem is an example.  Here, Ali's velocity at noon
      is $v(0) = -9$ m/min, and his velocity at 12:04 is $v(4) = 15$
      m/min; the average of these two is $\frac{-9 + 15}{2} = 3$ m/min,
      but his average velocity between noon and 12:04 was $-5$ m/min.

      If this seems counterintuitive, imagine a runner who takes the
      following hour-long run: he starts off at a furious pace of 10 mph,
      but he can only maintain that for 1 minute.  For the next 58
      minutes, he runs at 3 mph, and then he finds the motivation to run
      at 10 mph again for the last minute.  It should be intuitive that
      the runner's average speed for his entire run is close to 3 mph,
      since he ran at that speed for almost the entire hour.  But the
      average of his starting and ending speeds is 10 mph -- totally
      different than his average speed for the hour!
    \end{solution}

  \end{enumerate}

\pspace{\newpage}

\item 
  Suppose $F$ is an antiderivative of $f$ and $G$ is an antiderivative
  of $g$. Determine whether the following statements are true or false.
  \begin{enumerate}
  \item 
    \begin{problem}[\vfill]
      $FG$ is an antiderivative of $fg$.
    \end{problem}

    \begin{solution}
      \fbox{False}, and almost any example you try will show that this
      statement is false.  For example, $F(x) = x$ is an antiderivative of
      $f(x) = 1$ and $G(x) = 2x$ is an antiderivative of $g(x) = 2$, but
      $F(x) G(x) = 2x^2$ is not an antiderivative of $f(x) g(x) = 2$.
    \end{solution}

  \item 
    \begin{problem}[\vfill]
      $F + G$ is an antiderivative of $f + g$.
    \end{problem}

    \begin{solution}
      \fbox{True}: if $F' = f$ and $G' = g$, then $(F + G)' = f + g$. 
    \end{solution}

  \end{enumerate}

  \item Evaluate the following indefinite integrals.
  \begin{enumerate}[topsep=0pt]
    \item 
    \begin{problem}[\vfill]
        $\displaystyle \int\frac{2}{t}\,dt$
    \end{problem}
    
    \begin{solution}
        First, we can rewrite this as 
        $\displaystyle 2\int\frac{1}{t}\,dt$.
        We know that $\displaystyle \int \frac{1}{t}\,dt = \ln|t| + C$.
        
        Therefore $\displaystyle 2\int\frac{1}{t}\,dt =
        \boxed{2\ln|x| + C.}$
    \end{solution}
    
    \item 
    \begin{problem}[\vfill]
        $\displaystyle \int\frac{2}{t^2}\,dt$
    \end{problem}

    \begin{solution}
        First, we can rewrite this as 
        $\displaystyle 2\int\frac{1}{t^2}\,dt$.
        We know that $\displaystyle \int \frac{1}{t^2}\,dt = 
        -\frac{1}{t}+C$.
        
        Therefore $\displaystyle 2\int\frac{1}{t^2}\,dt =
        \boxed{-\frac{2}{t} + C.}$
    \end{solution}
    
    \item 
    \begin{problem}[\vfill]
        $\displaystyle \int\frac{2}{t^2+1}\,dt$
    \end{problem}

    \begin{solution}
        First, we can rewrite this as 
        $\displaystyle 2\int\frac{1}{t^2+1}\,dt$.
        We know that $\displaystyle \int \frac{1}{t^2+1}\,dt = 
        \arctan t +C$.
        
        Therefore $\displaystyle 2\int\frac{1}{t^2}\,dt =
        \boxed{2\arctan t + C.}$
    \end{solution}
  \end{enumerate}

  \item If time allows, try these using ``guess and check.''
  \begin{enumerate}[topsep=0pt]
    \item 
    \begin{problem}[\vfill]
        $\displaystyle \int\frac{1}{1+t}\,dt$
    \end{problem}

    \begin{solution}
        We begin by noticing that the integrand looks very similar to
        $\dfrac{1}{t}$. Since $\displaystyle \int \frac{1}{t}\,dt
        = \ln|t|+C$, we might reasonably guess that 
        $\displaystyle \int \frac{1}{1+t}\,dt
        \stackrel{?}{=} \ln|1+t|+C$.
        
        Differentiating to check:
        \begin{align*}
            \frac{d}{dx} \ln|1+t| &= \frac{1}{1+t} \cdot 1 \\
            & = \frac{1}{1+t}\, \checkmark
        \end{align*}
        Thus, $\displaystyle \int \frac{1}{1+t}\,dt
        = \boxed{\ln|1+t|+C.}$
    \end{solution}
    
    \item 
    \begin{problem}[\vfill]
        $\displaystyle \int\sin(t+2)\,dt$
    \end{problem}

    \begin{solution}
        We begin by noticing that the integrand looks very similar to
        $\sin(t)$. Since $\displaystyle \int \sin(t)\,dt
        = -\cos(t)+C$, we might reasonably guess that 
        $\displaystyle \int \sin(t+2)\,dt
        \stackrel{?}{=} -\cos(t+2)+C$.
        
        Differentiating to check:
        \begin{align*}
            \frac{d}{dx} -\cos(t+2) &= -(-\sin(t+2)) \cdot 1 \\
            & = \sin(t+1)\, \checkmark
        \end{align*}
        Thus, $\displaystyle \int \sin(t+2)\,dt
        = \boxed{-\cos(t+2)+C.}$
    \end{solution}
    
    \item
    \begin{problem}[\vfill]
        $\displaystyle \int\cos(3t)\,dt$
    \end{problem}

    \begin{solution}
        We begin by noticing that the integrand looks very similar to
        $\cos(t)$. Since $\displaystyle \int \cos(t)\,dt
        = \sin(t)+C$, we might reasonably guess that 
        $\displaystyle \int \cos(3t)\,dt
        \stackrel{?}{=} \sin(3t)+C$.
        
        Differentiating to check:
        \begin{align*}
            \frac{d}{dx} \sin(3t) &= \cos(3t) \cdot 3 \\
            & = 3\cos(3t)\
        \end{align*}
        This is not the original integrand. Our initial guess 
        $\sin(3t)$ was off by a factor of $3$. To fix this, we can
        divide our initial guess by $3$ to get 
        $\displaystyle \int \cos(3t)\,dt
        = \boxed{\tfrac{1}{3}\sin(3t)+C.}$
    \end{solution}
  \end{enumerate}
\end{enumerate}

\end{document}